% Implementation (WORK IN PROGRESS - Week 3)

% TODO: Expand this chapter after completing experiments
% TODO: Add code snippets and examples
% TODO: Create architecture diagrams
% TODO: Add screenshots of visualization

This chapter describes the implementation of the multi-agent fruit harvesting simulation. The system is built using Mesa 3.0, a Python framework for agent-based modeling.

\section{System Architecture}

% TODO: Add system architecture diagram here

The system consists of three main components:

\begin{itemize}
    \item \textbf{Environment}: 50×50 grid world with fruit resources
    \item \textbf{Agents}: Mobile harvester agents with communication capabilities
    \item \textbf{Data Collection}: Metrics tracking and visualization
\end{itemize}

\subsection{Overall Structure}

The implementation follows Mesa's standard architecture:

\begin{itemize}
    \item \textbf{Model class} (\texttt{HarvestModel}): Manages the simulation
    \item \textbf{Agent classes}: \texttt{HarvesterAgent} and \texttt{Fruit}
    \item \textbf{Grid}: \texttt{MultiGrid} for spatial representation
    \item \textbf{DataCollector}: Tracks metrics at each step
    \item \textbf{Visualization}: SolaraViz for interactive display
\end{itemize}

\subsection{Model Class}

% TODO: Add detailed explanation of HarvestModel class
% TODO: Include code snippet showing initialization

The \texttt{HarvestModel} class manages the simulation. Key parameters:

\begin{itemize}
    \item \texttt{width}, \texttt{height}: Grid dimensions
    \item \texttt{num\_agents}: Number of harvester agents
    \item \texttt{fruit\_density}: Proportion of cells with fruit
    \item \texttt{comm\_range}: Communication range in cells
    \item \texttt{seed}: Random seed for reproducibility
\end{itemize}

% TODO: Explain step() method in detail
% TODO: Add flowchart of simulation loop

\subsection{Grid and Spatial Representation}

% TODO: Add figure showing grid layout
% TODO: Explain boundary conditions

The environment is a 50×50 grid with non-toroidal boundaries. This creates realistic edge effects where agents near boundaries have fewer neighbors.

% TODO: Justify grid size choice
\end{itemize}

% TODO: Explain MultiGrid vs SingleGrid choice

\section{Agent Implementation}

% TODO: Add detailed agent architecture diagram

\subsection{Harvester Agent}

% TODO: Expand this section with detailed explanation

The \texttt{HarvesterAgent} class implements the harvesting behavior. Each agent maintains:

\begin{itemize}
    \item \texttt{pos}: Current position
    \item \texttt{harvested}: Count of fruit harvested
    \item \texttt{messages\_sent}: Communication counter
    \item \texttt{current\_target}: Target fruit location (if any)
    \item \texttt{visited\_cells}: Set of visited cells (for coverage)
\end{itemize}

\subsubsection{Agent Behavior}

% TODO: Add pseudocode or flowchart

The agent's \texttt{step()} method:
\begin{enumerate}
    \item Move (toward target or random walk)
    \item Harvest fruit if present
    \item Communicate with nearby agents
\end{enumerate}

% TODO: Explain movement algorithm in detail
% TODO: Explain harvesting logic
% TODO: Explain communication protocol

\textbf{Movement}: Agents move toward their target if they have one, otherwise they perform a random walk.

% TODO: Add code snippet showing movement logic

\textbf{Harvesting}: When an agent reaches a cell with fruit, it harvests the fruit and clears its target.

% TODO: Explain what happens to harvested fruit

\textbf{Communication}: Agents broadcast fruit locations to neighbors within communication range.

% TODO: Explain selective broadcasting (only to agents without targets)

\subsection{Fruit Agent}

% TODO: Expand fruit agent description

The \texttt{Fruit} class represents harvestable resources. Key properties:

\begin{itemize}
    \item \texttt{available}: Whether fruit can be harvested
    \item \texttt{times\_harvested}: Harvest counter
\end{itemize}

% TODO: Explain Static vs Replenishing dynamics (focus on Static for RQ1)

\section{Communication Protocol}

% TODO: Add detailed communication protocol diagram

\subsection{Range Calculation}

Communication range uses Chebyshev distance:

$$d_{Chebyshev}(a, b) = \max(|x_a - x_b|, |y_a - y_b|)$$

% TODO: Add figure showing Chebyshev distance examples
% TODO: Justify why Chebyshev vs Euclidean

\subsection{Message Passing}

% TODO: Expand this section with detailed explanation
% TODO: Add sequence diagram showing message flow

Messages contain fruit location coordinates. Agents broadcast to neighbors within range.

% TODO: Add detailed message processing algorithm
% TODO: Explain coordination mechanism

\subsection{Communication Cost}

% TODO: Expand cost model explanation

Communication cost is tracked by counting messages sent. Each message to a neighbor counts as one unit of cost.

% TODO: Justify this cost model vs alternatives (bandwidth, energy, etc.)

\section{Data Collection}

% TODO: Add detailed explanation of DataCollector setup

\subsection{Metrics Tracked}

The system tracks the following metrics at each step:

\begin{enumerate}
    \item \textbf{total\_yield}: Cumulative fruit harvested
    \item \textbf{cumulative\_messages}: Total messages sent
    \item \textbf{remaining\_fruit}: Available fruit on grid
\end{enumerate}

% TODO: Add coverage and efficiency metrics
% TODO: Explain how metrics are calculated
% TODO: Add code snippet showing DataCollector setup

\section{Parameter Selection}

% TODO: Expand this section with detailed justifications

\subsection{Grid Size: 50×50}

% TODO: Justify grid size choice with preliminary tests
% TODO: Compare with alternative sizes (20×20, 100×100)

Chosen to balance spatial dynamics with computational efficiency.

\subsection{Episode Length: 500 Steps}

% TODO: Run preliminary tests to determine optimal length
% TODO: Show convergence analysis

Preliminary tests suggest 500 steps is sufficient for agents to explore and harvest.

\subsection{Communication Ranges}

% TODO: Justify range values {0, 2, 4, 6, 8}
% TODO: Explain why these specific values

Testing ranges from 0 (no communication) to 8 (long range).

\subsection{Team Sizes and Densities}

% TODO: Justify team sizes {10, 20}
% TODO: Justify densities {0.15, 0.25}
% TODO: Explain interaction with grid size

Testing two team sizes (10, 20 agents) and two resource densities (0.15, 0.25).

\subsection{Replications}

% TODO: Justify 30 replications per configuration
% TODO: Discuss statistical power

% TODO: Calculate statistical power properly
% TODO: Justify 30 replications with power analysis

Planning to run 30 replications per configuration for statistical robustness.

\section{Testing and Validation}

% TODO: Write comprehensive test suite
% TODO: Add unit tests for each component

\subsection{Manual Testing}

% TODO: Document manual testing procedures

Currently testing manually by running simulations and checking:

\begin{itemize}
    \item Agents move correctly
    \item Fruit is harvested properly
    \item Messages are sent (need to verify counts)
    \item Metrics are collected
\end{itemize}

% TODO: Add automated tests
% TODO: Verify communication range calculation
% TODO: Check edge cases (agents at boundaries, etc.)

\subsection{Visualization Validation}

The Solara visualization helps debug agent behavior by showing:

\begin{itemize}
    \item Agent positions and movement
    \item Fruit locations
    \item Real-time metrics
\end{itemize}

% TODO: Add communication links visualization
% TODO: Improve visualization performance

\section{Experimental Infrastructure}

% TODO: Write batch experiment runner script
% TODO: Implement progress tracking
% TODO: Add result saving to CSV

\subsection{Planned Batch Runner}

Need to write a script to automate the 600 experimental runs:

\begin{itemize}
    \item Generate all 20 configurations
    \item Run 30 replications per configuration
    \item Save results to CSV
    \item Track progress
\end{itemize}

% TODO: Implement this script (Week 3-4)

\subsection{Reproducibility}

Using fixed random seeds (0-29) for reproducibility.

% TODO: Verify seed handling in Mesa 3.0
% TODO: Test that same seed produces same results

\section{Implementation Challenges}

% TODO: Document challenges encountered and solutions

\subsection{Mesa 3.0 API}

Learning Mesa 3.0 API (different from 2.x):

\begin{itemize}
    \item No \texttt{unique\_id} in agent constructors
    \item Agent activation via \texttt{shuffle\_do('step')}
    \item Visualization requires model instance
\end{itemize}

% TODO: Add more details about API differences

\subsection{Known Issues}

% TODO: Fix these issues before running full experiments

Current issues to resolve:

\begin{itemize}
    \item Communication range calculation needs verification
    \item Agents sometimes cluster in corners
    \item Visualization is slow with many agents
    \item Need to add coverage and efficiency metrics
\end{itemize}

\section{Summary}

% TODO: Expand summary after completing implementation

This chapter outlined the implementation approach. The system uses Mesa 3.0 with a grid-based environment, harvester agents, and communication protocol. Basic functionality is working, but several components need completion before running full experiments.

% TODO: Write Results chapter after experiments complete

